\subsection{Traditional Approaches}\label{subsec:traditional}

As traditional approaches, we consider the three established local adaptive thresholding techniques: Niblack~\cite{Niblack1985}, Sauvola~\cite{Sauvola1997}, and Nick~\cite{Khurshid2009}.
All three approaches compute a pixel-wise threshold based on statistics extracted from a local window centered around each pixel.
The binarized output is then obtained by comparing the pixel intensity against the locally computed threshold. 

% Niblack
\paragraph{Niblack Thresholding.}
Niblack's method computes the local mean intensity $m$ and standard deviation $\sigma$ within a sliding window.
The threshold is defined as

\begin{equation}
    T = m + k \cdot \sigma ,
\end{equation}

where $k$ is a user-defined parameter controlling the influence of local contrast. 
In our experiments, we use $k=-0.2$, as suggested in the original formulation.
Each pixel is classified according to whether its grayscale intensity is below or above the threshold, producing a binary image.

% Sauvola
\paragraph{Sauvola Thresholding.}
Sauvola and Pietikäinen~\cite{Sauvola1997} extend Niblack's approach by introducing a normalization term that improves robustness in low-contrast and noisy regions.
The thresholding is computed as

\begin{equation}
    T = m \left(1 + k \left(\frac{\sigma}{R} - 1\right)\right) ,
\end{equation}

where $R$ denotes the dynamic range of the standard deviation (commonly set to $R=128$ for 8-bit images).
The parameter $k$ typically lies in the range $[0.2, 0.5]$. 
In our implementation, we select $k = 0.34$ and $R = 128$. 

% Nick
\paragraph{Nick Thresholding.}
Khurshid et al.~\cite{Khurshid2009} propose the Nick method to address contrast degradation observed in Sauvola's formulation.
Instead of directly using $\sigma$, the method incorporates the second moment of the local intensity distribution:

\begin{equation}
    T = m + k \cdot \sqrt{\frac{(\sum P_i^2 - m^2)}{N}},
\end{equation}

where $P_i$ denotes the pixel intensities within the local window and $N$ is the number of pixels in that window.
The parameter $k$ is typically chosen in the range $[-0.2, -0.1]$.
In our experiments, we use $k = -0.15$.
